% Targeted replacements and additions for Chapter 4 (Methodology).
% These blocks implement comments on interdisciplinarity, traceability,
% manual review, and the special role of the VADA corpus.

% =============================================================================
% 1. Replace Section 4.1, ``Overview of the approach''.
% =============================================================================

\section{Overview of the approach}

The methodology follows a document-processing pipeline designed to transform
heterogeneous textual resources into a structured corpus that can support search,
filtering, comparison, and thematic analysis. Its objective is to make dispersed
descriptions of territorial initiatives related to population ageing more accessible
for research, participatory action, and future operational use.

The methodology was developed through an interdisciplinary dialogue between the
technical development of the Index and the social-science researchers of the SIAGE
Chair. This dialogue informed the selection of sources, the definition of the
metadata schema, the vocabulary used for relevance scoring, and the interpretation
of the resulting data. The pipeline should therefore not be understood as a purely
technical procedure applied to a pre-existing corpus. It is a means of
capitalising, comparing, and circulating knowledge produced both through Senior
Citizen Design Workshops and through initiatives documented by other territorial
and gerontological actors.

The methodological process evolved during the internship. The initial approach
relied on a multi-step rule-based pipeline: documents were cleaned, segmented,
analysed for document sections, and processed to extract candidate metadata
values. Candidate values are potential values for a given metadata field, such as
a place name that may describe a territory or an organisation that may be an
initiative holder. These candidates were scored according to explicit rules. This
approach was transparent and useful for inspecting intermediate results, but the
diversity of source formats and writing styles made it difficult to generalise the
extraction rules to every document type.

The final approach therefore complements the rule-based pipeline with LLM-assisted
information extraction. A local language model receives a constrained prompt and
transforms each textual document into a structured JSON record. The methodology
is organised around six components: collection, corpus construction, cleaning and
quality control, progressive rule-based structuring, constrained LLM extraction,
and human evaluation.

% =============================================================================
% 2. Add this paragraph at the END of ``Data collection''.
% =============================================================================

Source-specific collection was necessary not only because formats differed, but
also because relevant information was distributed unevenly within and between
websites. In some sources, an initiative was described on a single page; in others,
its title, objectives, territorial scope, and actors were spread across several pages
or embedded in a PDF. The collection scripts were therefore designed to navigate
the editorial structure of each source and to retain the text needed to understand
an initiative in context. This work was carried out with attention to the
documentary practices of the selected actors, rather than assuming that all sources
could be treated through a single generic scraping procedure.

% =============================================================================
% 3. Replace the FIRST TWO paragraphs of ``Corpus construction'' with this text.
% Keep the existing table and the following paragraphs after it.
% =============================================================================

The corpus was built from two complementary forms of documentary material. First,
the SIAGE Chair provided action sheets produced by older adults themselves within
Senior Citizen Design Workshops. These documents record locally identified issues,
proposed responses, target publics, and implementation conditions. Second,
external institutional, associative, and territorial sources were selected because
they document initiatives, programmes, research projects, or public-action
resources relevant to the adaptation of territories to population ageing.

The selection of external sources was conducted in dialogue with the
social-science researchers of the SIAGE Chair. The objective was to include actors
whose material was both relevant to the gerontological field and sufficiently
substantial to support documentary structuring. Only documents containing a
meaningful description of an initiative, programme, action, or research project
were retained. Navigation pages, empty pages, duplicated pages, and pages without
a clear connection to territorial adaptation to ageing were excluded automatically
when possible and otherwise reviewed manually.

% =============================================================================
% 4. Add this paragraph AFTER Table 4.1 (final corpus).
% =============================================================================

The sources differ in institutional role and documentary purpose. CNAV documents
prevention-related initiatives and institutional resources; the Fondation de France
publishes accounts of associative and territorial initiatives; HCFEA provides
public-policy and institutional resources; IReSP Autonomie documents research
projects related to autonomy; and VADA records local age-friendly initiatives.
These sources were not selected as statistically equivalent samples. They were
selected because, together with the SIAGE action sheets, they make visible several
forms of knowledge and action related to territorial adaptation to ageing.

% =============================================================================
% 5. Replace the paragraph beginning ``The corpus is not a statistically
% representative sample...'' with this text.
% =============================================================================

The corpus is not a statistically representative sample of all French initiatives
related to ageing. VADA contributes 854 of the 1,153 records, or approximately
74\% of the corpus, and this predominance must be considered when interpreting
descriptive statistics. At the same time, VADA is not merely an overrepresented
source: the Francophone Age-Friendly Cities and Communities network plays a
structuring role in identifying, documenting, and disseminating local initiatives
developed by participating territories. Its prominence therefore reflects both a
source imbalance and its role as a national resource for the capitalisation of local
ageing-policy experience. Results are consequently interpreted as descriptions of
the collected documentation, not as estimates of the complete French landscape.

% =============================================================================
% 6. Add this paragraph BEFORE Table 4.2 (keyword weighting).
% =============================================================================

The keywords and their weights were not defined arbitrarily. They were developed
through discussion with the SIAGE team in order to reflect the scientific scope of
the Index. Concepts that are central to population ageing and territorial adaptation
were given higher weights, while broader terms associated with relevant publics or
contexts were retained with lower weights. The resulting vocabulary is therefore a
practical quality-control instrument informed by gerontological expertise, rather
than a claim to provide an exhaustive taxonomy of the field.

% =============================================================================
% 7. Replace the two paragraphs beginning ``For each document...'' and ``This
% relevance score is therefore...'' with this text.
% =============================================================================

For each document, the script counts the occurrences of these keywords in the
normalised text and applies a small bonus when a keyword appears in the document
title or file name. The resulting score is normalised by the number of words in the
document to make documents of different lengths more comparable. During the
project, ranked lists of high- and low-scoring documents were reviewed with the
SIAGE team. This review made it possible to examine extreme cases in the light of
gerontological expertise and to decide whether a document should be retained,
excluded, or inspected further.

The relevance score is therefore a quality-control and manual-review
prioritisation mechanism. It is not a supervised classifier, a validated measure of
relevance, or an evaluation of the factual quality of subsequent metadata
extraction.

% =============================================================================
% 8. Replace Section 4.5, ``Progressive structuring of the document database''.
% =============================================================================

\section{Progressive structuring of the document database}

The rule-based architecture uses separate directories for each major
transformation. These intermediate databases are not independent corpora: they are
successive, inspectable representations of the same documents. Their purpose is to
make it possible to identify at which stage an extraction error or an inappropriate
candidate value was introduced, and to rerun a stage after modifying a rule without
overwriting the original material.

\begin{verbatim}
Index_Insight/
|-- Database/
|-- Cleaned_Database/
|-- Segmented_Database/
|-- Sectioned_Database/
|-- Entities_Database/
|-- Candidates_Database/
|-- Scored_Database/
|-- Final_Tagged_Database/
|-- LLM_Tagged_Database/
|-- Evaluation/
|-- outputs_graphs/
|-- Scripts/
\end{verbatim}

The \texttt{Database} directory contains the raw collected text files and
\texttt{Cleaned\_Database} contains their normalised versions. The
\texttt{Segmented\_Database} stores paragraph-, line-, and sentence-level
representations. The \texttt{Sectioned\_Database} records detected headings or
document sections, such as objectives, target publics, context, stakeholders, or
territorial information, when they are explicitly present.

The \texttt{Entities\_Database} stores intermediate entity-extraction results.
The \texttt{Candidates\_Database} then gathers potential values for each metadata
field, while \texttt{Scored\_Database} associates them with heuristic scores and
confidence indicators. Finally, \texttt{Final\_Tagged\_Database} consolidates the
selected values into a rule-based JSON record and produces a global JSON file for
review and analysis.

This organisation supported iterative development. For example, when manual
inspection revealed recurrent ambiguities in territorial scales or territorial names,
the associated rules and controlled values could be revised and the relevant stages
rerun. A final consolidation stage then brought the successive outputs together in a
cleaned, machine-readable dataset. This traceability is an advantage of the
rule-based baseline, even though manually designed rules do not generalise easily
across the full diversity of documents.

% Add after the sentence introducing the expected JSON schema, before its list.
The schema was co-constructed with the SIAGE team. Its fields are not intended
only to structure information technically; they were selected because they support
the methodological needs of the Senior Citizen Design approach. For example,
\texttt{problematique}, \texttt{objectifs}, and \texttt{solution\_envisagee}
support the description of local issues and responses; \texttt{territoire} and
\texttt{echelle} support territorial comparison; and \texttt{porteur\_initiative},
\texttt{public\_vise}, \texttt{thematique}, and \texttt{nature\_initiative}
support the comparison of actors, publics, themes, and intervention formats. The
schema thereby supports diagnosis, co-design, dialogue with decision-makers, and
the capitalisation of experience, while improving the homogeneity and comparability
of records between territories.

% Add after the paragraph on controlled vocabularies in the LLM section.
For interpretability, the thematic labels include examples such as health and
prevention, mobility, housing, social connection, digital inclusion,
intergenerational relations, and citizen participation. Initiative types include, for
example, workshops, training, awareness-raising activities, services, support
schemes, digital tools, events, publications, dedicated places, and planning
projects. These labels are analytical categories: they make it possible to organise
the corpus while retaining the source text and provenance needed to interpret a
record in context.

% Add before the paragraph beginning ``The current configuration uses...''.
VADA required particular attention because it is the largest source in the corpus
and because its initiative descriptions are internally varied in length, style, and
level of detail. The dedicated VADA extraction script was therefore used to test
prompt constraints and input handling on a large, heterogeneous set of local
initiative descriptions. Its configuration should not be interpreted as a separate
research question, but as a source-specific implementation of the common metadata
schema.

% Add after the paragraph beginning ``Raw model responses are archived...''.
Validation was conducted at two complementary levels. During iterative
development, an initial random sample of ten records was reviewed manually with
the SIAGE team in order to identify recurring errors, inconsistencies, and
unsatisfactory field values. The resulting observations informed corrections to
rules, controlled vocabularies, and prompt instructions. This exploratory review is
distinct from the formal human reference evaluation reported in Chapter 6, which
uses a stratified sample of 60 documents and field-level assessment criteria.
